\section{Discussion}\label{sec:discussion}
\subsection{Interpretation and comparison with earlier studies}
\textbf{RQ1.} Headline sentiment in large caps mostly confirms price moves that have already happened, which fits Tetlock's temporary price pressure~\cite{tetlock2007} and the over-reaction to heavily covered news reported by Hillert et al.~\cite{hillert} better than the slow absorption of weekly tone found by Heston and Sinha~\cite{hestonsinha}. The difference is consistent with Ke et al.~\cite{kekellyxiu}, who find that news about large firms is priced within days: the firms here are heavily covered, and dates-only news reaches the portfolio a day or more after professional feeds.

\textbf{RQ2.} Like Coriat and Benhamou~\cite{coriat}, the allocation that combined price and sentiment did better than the price-only allocation, with a higher return and Sharpe ratio and a smaller drawdown. Unlike in their study, the correction did not make the forecast more accurate, and the gain is not statistically significant. The gain is still economically relevant. In liquid large caps a credible edge is expected to be small and spread over many decisions~\cite{grinoldkahn}, and an extra 1--2 percentage points a year, earned on both buy and sell decisions, matters to an investor, even if eleven years are too short to separate it from luck.

\textbf{RQ3.} Using sentiment as Black--Litterman views, as Colasanto et al.\ did~\cite{colasanto}, gave a marginally higher Sharpe ratio and a smaller drawdown than the benchmarks, with little trading. DeMiguel et al.'s warning that $1/N$ is hard to beat~\cite{demiguel} still holds: no difference is statistically significant over eleven years. In economic terms, a drawdown about five percentage points smaller than equal weighting at a similar return is valuable to an investor who must hold through sell-offs.

\subsection{Implications for portfolio managers}
The results give two guidelines. First, public headline sentiment is better used together with a price-based forecast than as a stand-alone source of alpha in liquid large caps; here it improved the allocation without improving forecast accuracy, and gains of this size justify further testing even though one sample cannot confirm them. Second, any signal should enter through a skill-weighted, cost-aware optimiser, which limits the damage when a signal stops working. Deployment is a weekly batch job of headline retrieval, FinBERT scoring and one quadratic programme, which runs in minutes on a laptop.

\subsection{Limitations and future work}
The universe excludes nine firms that later disappeared, a bias of unknown sign. The news archive was retrieved in 2026 and ranked by today's relevance, which may favour items still online and stories that later proved important. Many comparisons were made, eleven years give little power to detect Sharpe differences of about 0.1, the cost penalty was not tuned, and MSCI World also contains markets outside the universe. Weekly decisions cannot capture effects that last only hours. Future work should use a real-time news feed with time stamps, a broader mid-cap universe where news is absorbed more slowly, and a cost penalty and holding horizon chosen on the validation period.
